\chapter*{Introduction Générale\label{intro}} 
\markboth{\uppercase{Introduction Générale}}{\uppercase{Introduction Générale}}
\addcontentsline{toc}{chapter}{Introduction Générale}

Le passage à l'économie numérique et la généralisation des services en ligne en Algérie ont entraîné une production massive de données à caractère personnel. Dans ce contexte, la protection de la vie privée est devenue une exigence constitutionnelle et légale majeure. La promulgation de la Loi n° 18-07 du 10 juin 2018, relative à la protection des personnes physiques dans le traitement des données à caractère personnel, a marqué un tournant décisif pour les organisations nationales.

Cependant, face à l'accélération des cybermenaces et à l'émergence de l'intelligence artificielle, le législateur a renforcé ce dispositif par la Loi n° 25-11 relative à la sécurité et à la souveraineté numérique. Ce nouveau texte impose désormais une maîtrise totale des infrastructures de stockage et une régulation stricte des algorithmes de traitement, consolidant ainsi le principe de localisation des données sur le territoire national.

Désormais, toute entité traitant des données (qu'il s'agisse d'un organisme public ou d'une entreprise privée) est investie de responsabilités juridiques et techniques strictes, régies par la complémentarité de ces deux lois. Le non-respect de ces dispositions expose les responsables de traitement à des sanctions administratives et pénales lourdes, supervisées par l'Autorité Nationale de Protection des Données à Caractère Personnel (ANPDP).

Toutefois, la mise en conformité ne se limite plus à un simple aspect juridique ; elle nécessite une refonte profonde de l'architecture des systèmes d'information pour garantir une souveraineté technologique. C'est ici que réside l'enjeu de notre étude : comment transformer ces contraintes légales en une architecture logicielle automatisée, sécurisée par un chiffrement de bout en bout et assistée par une Intelligence Artificielle souveraine (LLM Local) ?

Ce premier chapitre s'attache à définir le cadre conceptuel de notre projet. Nous y aborderons les principes fondamentaux des lois 18-07 et 25-11, les rôles des différents acteurs impliqués, ainsi qu'une analyse de l'état de l'art concernant les technologies de "Compliance by Design".